\chapter{Methodology}
\label{ch:methodology}

\section{System Architecture}
\label{sec:method_arch}

% TODO: Block diagram (figure) of the full system:
%   mic -> ADC -> nn~ (BRAVE) -> gain -> DAC -> speakers
% Describe each block briefly.
% State the design constraint: sub-10ms end-to-end latency.
% Explain how causal convolutions in BRAVE enable streaming.

\subsection{BRAVE in one paragraph}
% TODO: Streaming-RAVE extension; causal padding; PQMF multi-band
% decomposition; multi-scale STFT reconstruction loss; latent VAE
% with KL regularisation; Phase 2 GAN refinement via multi-scale
% discriminator.

\subsection{Real-Time Deployment Constraints}
% TODO: Why Pure Data + nn~. The nn~ external loads TorchScript-exported
% models and runs forward() per audio buffer. Block size 64 samples at
% 44.1kHz = 1.45ms processing window. Causal-only architecture (no
% lookahead) keeps total latency below 10ms.

\section{Training Paradigm}
\label{sec:method_training}

\subsection{Two-Phase Training}
% TODO: Phase 1 = reconstruction (encoder + decoder, MS-STFT loss + KL).
% Phase 2 = adversarial refinement (decoder vs discriminator).
% Explain why this works: stable reconstruction first, then perceptual polish.

\subsection{Beta Schedule and Posterior Collapse}
% TODO: KL weight beta controls regularisation vs reconstruction trade-off.
% Constant beta from step 0 collapses the posterior (guitar_v1 failure).
% Log-space warmup from 1e-4 to target over Phase 1 prevents this.
% Reference BRAVE paper / our v1 vs v2 comparison.

\subsection{Latent Size Selection}
% TODO: BRAVE default = 128. Our v2 with 128 produced 0.005 nats/dim
% (effectively unused). IIL pretrained guitar uses 16. We use 16
% (proven empirically in v3-v6).

\section{Datasets and Preprocessing}
\label{sec:method_data}

\subsection{Guitar Corpus}
% TODO: Table of sources, durations, recording methods, totals.
% GuitarSet (3.05h mic'd), Guitar-TECHS (~2h DI mono after filtering
% out stereo files), IDMT-SMT-Guitar dataset4 (4.35h, mixed mic/pickup/DI),
% IDMT-SMT-Guitar dataset2 (1.02h, removed due to 24-bit incompatibility).
% Final: ~9.4h after cleanup. LMDB n_seconds = 57665 (data augmented
% via overlapping windows during preprocessing).

\subsection{Drums Corpus}
% TODO: Groove MIDI Dataset audio. 1090 files, ~10.86h after stereo-to-mono
% conversion. Acoustic drums (vs BRAVE paper's electronic drums).

\subsection{Preprocessing Pipeline}
% TODO: Steps:
%   1. Walk source directories, find WAV files.
%   2. Validate (librosa.load); skip empty or corrupt files.
%   3. Stereo -> mono via librosa mono=True.
%   4. Resample to 44.1kHz if needed.
%   5. Save as PCM_16 to consolidated directory.
%   6. rave preprocess --no-lazy to build LMDB.
% Note the failure modes discovered along the way:
%   - 24-bit audio files silently deadlock rave preprocess.
%   - Stereo files do the same.
%   - __MACOSX/ metadata .wav files (empty) hang rave preprocess.
% These are documented as practical findings.

\subsection{Noise-Floor Analysis}
\label{sec:noise_floor}
% TODO: Reproduce the noise-floor table from setup_notes.md.
% Method: 5th percentile of 100ms RMS windows, per source.
% Result: 35.8 dB spread across recording methods.
% Implication: encoder will encode recording method as a latent feature.

\section{Diagnostic Methodology}
\label{sec:method_diagnostics}

\subsection{TensorBoard Metrics}
% TODO: What each metric means and what healthy looks like.
% regularization (KL): rises through Phase 1, plateaus in Phase 2.
% multiband_spectral_distance: drops through Phase 1, briefly spikes
%   at Phase 2 onset (GAN disrupts reconstruction), then settles.
% pred_real / pred_fake: discriminator confidence. Gap < 3 = healthy,
%   gap > 5 = discriminator dominance (failure indicator).
% loss_dis, loss_gen: balance shows GAN health.

\subsection{Python Diagnostic Test}
% TODO: We pass five distinct inputs through the trained model and
% compare encoder latents and decoder outputs:
%   - silence
%   - sine 220 Hz, sine 440 Hz, sine 880 Hz
%   - frequency sweep 100 - 2000 Hz
%   - random noise
%
% Metrics:
%   - Encoder differentiation: pairwise mean absolute difference of latents.
%   - Decoder responsiveness: pairwise mean absolute difference of outputs.
%   - Output amplitude (signal vs silence): a healthy model produces less
%     output for silence than for signal.
%
% This test takes about 30 seconds per checkpoint and is run after each
% training iteration.

\subsection{Reconstruction Test}
% TODO: We pass real training-distribution audio (guitar from GuitarSet)
% through the model and compare input vs output amplitude and STFT
% magnitude correlation. A working autoencoder should retain amplitude
% and show STFT correlation > 0.3.
